% =====================================================================
%  Development Environment and Toolchain
%  SCOPE: Exact versions, installation, build and launch procedures.
% =====================================================================
\section{Development Environment and Toolchain}
\label{sec:environment}

This section specifies the software environment precisely enough that
every result in this report can be regenerated on a clean machine. It is
placed before the technical chapters deliberately: several of the
defects analysed later are consequences of version-specific behaviour
(a deprecated Gazebo service, a build mode that resolves entry points
through a symbolic link), and those are not intelligible without knowing
exactly what was running.

\subsection{Hardware and operating system}

All development and all reported experiments were carried out on a
single workstation, Table~\ref{tab:hardware}. No GPU-accelerated
inference is used anywhere in the stack; the only GPU consumer is the
simulator's rendering pipeline, which also serves the simulated lidar
(Sec.~\ref{sec:phaseB_twin} explains why that matters).

\begin{table}[!t]
\caption{Development and experiment platform}
\label{tab:hardware}
\centering\footnotesize
\begin{tabular}{@{}ll@{}}
\toprule
\textbf{Item} & \textbf{Value}\\
\midrule
Machine            & HP EliteBook laptop\\
Operating system   & Ubuntu 24.04 LTS (Noble Numbat)\\
Display server     & Wayland, with Qt forced to XCB\\
Python             & 3.12 (system interpreter)\\
Project storage    & external volume, mounted under \filename{/mnt}\\
\bottomrule
\end{tabular}
\end{table}

The display server is worth recording. Gazebo's GUI detects Wayland and
falls back to the XCB Qt platform plugin, printing
\texttt{Detected Wayland. Setting Qt to use the xcb plugin} at every
start. This is benign in itself, but it is one reason the GUI
camera-tracking behaviour differs from the documented default, which
motivated the camera watchdog described in Sec.~\ref{sec:phaseB_twin}.

\subsection{Software versions}

Table~\ref{tab:versions} lists every component whose version affects the
results. Two entries are load-bearing:

\begin{itemize}[leftmargin=*,itemsep=1pt,topsep=2pt]
  \item \textbf{Gazebo Harmonic (gz-sim 8.11.0)} deprecated the
        \topicname{gui/follow} service in favour of the
        \topicname{gui/track} topic. Code written against the older API
        fails silently --- the camera simply never locks onto the robot
        --- which cost one full debugging cycle
        (Appendix~\ref{app:troubleshooting}).
  \item \textbf{\texttt{colcon build -{}-symlink-install}} resolves
        console-script entry points through an egg-link back into
        \filename{src/}. When that metadata goes missing, every node
        dies at import with \texttt{PackageNotFoundError} --- fifteen
        identical tracebacks that name the symptom and not the cause.
        The launch wrapper tests for this explicitly before starting
        anything.
\end{itemize}

\begin{table}[!t]
\caption{Software components and versions}
\label{tab:versions}
\centering\footnotesize
\begin{tabular}{@{}lll@{}}
\toprule
\textbf{Component} & \textbf{Version} & \textbf{Role}\\
\midrule
Webots                        & R2025a          & \phaseA{} simulator\\
ROS~2                         & Jazzy Jalisco   & \phaseB{} middleware\\
Gazebo Sim                    & Harmonic 8.11.0 & \phaseB{} simulator\\
\texttt{ros\_gz\_sim}         & Jazzy           & spawn, world launch\\
\texttt{ros\_gz\_bridge}      & Jazzy           & topic bridging\\
\texttt{webots\_ros2\_driver} & Jazzy           & \phaseA{}$\rightarrow$ROS\\
RViz~2                        & Jazzy           & visualisation\\
NumPy                         & system package  & grids, scan matching\\
colcon                        & system package  & build tool\\
Physics engine                & DART (default)  & 1\,ms step\\
\bottomrule
\end{tabular}
\end{table}

\subsection{Installation}

Listing~\ref{lst:install} gives the complete installation from a clean
Ubuntu 24.04. The ROS~2 and Gazebo repositories are added first; the
\texttt{ros-jazzy-ros-gz} metapackage pulls in both the bridge and the
simulator integration.

\begin{lstlisting}[style=bash,caption={Full installation on a clean
Ubuntu 24.04.},label={lst:install}]
# --- 1. ROS 2 Jazzy -------------------------------------
sudo apt update && sudo apt install -y \
     software-properties-common curl
sudo add-apt-repository universe
sudo curl -sSL -o /usr/share/keyrings/ros-archive-keyring.gpg \
  https://raw.githubusercontent.com/ros/rosdistro/master/ros.key
echo "deb [signed-by=/usr/share/keyrings/\
ros-archive-keyring.gpg] \
http://packages.ros.org/ros2/ubuntu noble main" \
     | sudo tee /etc/apt/sources.list.d/ros2.list
sudo apt update
sudo apt install -y ros-jazzy-desktop

# --- 2. Gazebo Harmonic + the ROS 2 bridge --------------
sudo apt install -y ros-jazzy-ros-gz \
                    ros-jazzy-ros-gz-sim \
                    ros-jazzy-ros-gz-bridge

# --- 3. Webots R2025a (Phase A) -------------------------
sudo apt install -y ros-jazzy-webots-ros2
# plus the Webots R2025a .deb from cyberbotics.com

# --- 4. Build tooling and Python dependencies -----------
sudo apt install -y python3-colcon-common-extensions \
                    python3-numpy python3-opencv

# --- 5. The report toolchain (optional) -----------------
sudo apt install -y texlive-latex-extra texlive-publishers \
                    texlive-science
\end{lstlisting}

\subsection{Workspace layout and build}

The ROS~2 workspace contains five packages, Table~\ref{tab:packages}.
The split is not cosmetic: \texttt{youbot\_control} holds the algorithms
and the control nodes and depends on \emph{no} simulator, while
\texttt{youbot\_gazebo} and \texttt{youbot\_webots} are thin,
simulator-specific shells. This is what allowed \phaseA{} to be reused
in \phaseB{} rather than rewritten (Sec.~\ref{sec:architecture}).

\begin{table}[!t]
\caption{ROS~2 workspace packages and their dependencies}
\label{tab:packages}
\centering\footnotesize
\begin{tabular}{@{}p{0.24\columnwidth}p{0.66\columnwidth}@{}}
\toprule
\textbf{Package} & \textbf{Depends on}\\
\midrule
\texttt{youbot\_control} &
  \texttt{rclpy}, \texttt{geometry\_msgs}, \texttt{sensor\_msgs},
  \texttt{nav\_msgs}, \texttt{tf2\_ros},
  \texttt{visualization\_msgs}, \texttt{numpy}\\
\texttt{youbot\_slam} &
  \texttt{rclpy}, \texttt{nav\_msgs}, \texttt{sensor\_msgs},
  \texttt{geometry\_msgs}, \texttt{tf2\_ros},
  \texttt{youbot\_control}\\
\texttt{youbot\_gazebo} &
  \texttt{youbot\_control}, \texttt{youbot\_bringup},
  \texttt{ros\_gz\_sim}, \texttt{ros\_gz\_bridge},
  \texttt{robot\_state\_publisher}, \texttt{rviz2}\\
\texttt{youbot\_bringup} &
  \texttt{youbot\_control}, \texttt{launch}, \texttt{launch\_ros},
  \texttt{rviz2}, \texttt{robot\_state\_publisher}\\
\texttt{youbot\_webots} &
  \texttt{webots\_ros2\_driver}, \texttt{youbot\_control},
  \texttt{rclpy}\\
\bottomrule
\end{tabular}
\end{table}

Note the direction of the arrows: no package containing an algorithm
depends on a simulator, and no simulator package is depended upon by
another. \texttt{youbot\_slam} depends on \texttt{youbot\_control} only
to reuse the occupancy-grid and scan-matching libraries.

\begin{lstlisting}[style=bash,caption={Cloning and building.},
label={lst:build}]
git clone https://github.com/\
franckarmandjosiasEinstein-bit/ROS2_stage.git
cd ROS2_stage
source /opt/ros/jazzy/setup.bash
colcon build --symlink-install
source install/setup.bash
\end{lstlisting}

\texttt{-{}-symlink-install} is used so that editing a Python node takes
effect without a rebuild, which matters when one simulation run costs
half an hour of wall time. Its price is the entry-point fragility
described above; if that appears, a plain \texttt{colcon build} writes
real distribution metadata and removes the dependency on
\filename{src/} entirely.

\subsection{Running a simulation}

Every experiment is launched through one wrapper script,
\filename{src/youbot\_gazebo/scripts/run.sh}, rather than through
\texttt{ros2 launch} directly. The wrapper exists because of three
concrete failures, each of which cost at least one run:

\begin{enumerate}[leftmargin=*,itemsep=1pt,topsep=2pt]
  \item \textbf{Gazebo does not die on Ctrl\,+\,C.} A surviving
        \texttt{gz sim} keeps publishing an old \topicname{clock}; the
        next run then has two clocks, RViz thrashes, and the GUI can
        attach to a ghost world with no robot in it. The wrapper
        installs a cleanup trap and afterwards \emph{verifies} that
        nothing survived, printing the command line of any process it
        could not match rather than leaving it to be inferred from the
        next run misbehaving.
  \item \textbf{Gazebo rewrites its GUI configuration on every exit}
        with wherever the camera was left, and reads it back in
        preference to the world's \texttt{<gui>} pose. A run started
        after the camera was parked on empty space opens blind. The
        wrapper deletes \filename{\textasciitilde/.gz/sim/8/gui.config}
        before starting.
  \item \textbf{Silent build breakage}, via the entry-point check
        described above.
\end{enumerate}

The wrapper then runs the pre-flight regression suite
(Sec.~\ref{sec:regression}) and refuses to start if any check fails:
about one second of pure Python weighed against thirty minutes of
simulation.

\begin{lstlisting}[style=bash,caption={The supported launch
configurations.},label={lst:run}]
# Level 2 -- ground-truth pose, no SLAM (baseline)
./src/youbot_gazebo/scripts/run.sh

# Level 3 -- homemade SLAM in the loop
./src/youbot_gazebo/scripts/run.sh slam

# Level 3 -- calibrated dead reckoning, SLAM disabled
#            (the production configuration, Sec. V-E)
./src/youbot_gazebo/scripts/run.sh slam \
      slam:=false pose_topic:=odom_calibrated

# Headless, for batch runs
./src/youbot_gazebo/scripts/run.sh slam gui:=false
\end{lstlisting}

A correct start prints three markers, in this order. Their absence is
the fastest way to detect that a stale build is being run --- a mistake
made twice during this project, both times producing a log that looked
plausible and was measuring the wrong binary:

\begin{lstlisting}[style=bash]
[run.sh] pre-flight: all 191 checks passed.
[mapping_node] mapping_node up: /scan + /odom -> /map
               (inflated, for planning) + /map_raw (evidence)
[perf_monitor] perf_monitor up: timing, throughput and
               where the time goes.
\end{lstlisting}

\subsection{Running \phaseA{} (Webots)}

\phaseA{} is self-contained and does not require ROS. Webots is opened
on one of the two worlds and the \texttt{controller} field of the YouBot
node selects which of the fifteen controllers runs
(Sec.~\ref{sec:phaseA}). The \texttt{youbot\_webots} package exists only
to expose the same robot to ROS~2 through
\texttt{webots\_ros2\_driver}; it was used to confirm that the algorithm
layer behaved identically under both simulators before the Gazebo twin
was built.

\begin{lstlisting}[style=bash,caption={Launching \phaseA{}.},
label={lst:webots}]
# Standalone Webots (no ROS): open the world, then choose
# the controller in the scene tree.
webots worlds/smart_agriculture.wbt
webots worlds/greenhouse.wbt

# The same robot exposed to ROS 2 (bridge validation)
ros2 launch youbot_webots webots.launch.py
\end{lstlisting}

\subsection{Version control and reproducibility}

The two phases live in two repositories: the Webots prototype in
\texttt{Projet}, and the ROS~2 stack in \texttt{ROS2\_stage}. Every
result quoted in Sec.~\ref{sec:results} names the commit it was produced
from, because two of the analyses in this report compare runs that
differ by a single parameter file. Appendix~\ref{app:reproduce} lists
those commits alongside the exact command and the expected output.
